\chapter{Statistical Analysis}
% \label{chap:statistical_analysis}

\section{Analytical Aim and Data Provenance}
This chapter analyses the empirical outputs produced by the experimental system described in
\Cref{chap:arch_design} and the design factors in \Cref{tab:design_factors}. The purpose is not
only to rank ASR models, but to explain which experimental variables account for changes in
transcription error and where the comparison is statistically well controlled.

The analysis uses the current top-level result files:
\texttt{WER\_results\_faster-whisper.json},
\texttt{WER\_results\_whisperx.json},
\texttt{WER\_results\_wav2vec2.json}, and
\texttt{WER\_results\_parakeet.json}. The accompanying figures and the numerical summary file
\texttt{Documentation/statistical\_analysis\_summary.json} were generated from these files using
\texttt{Code/generate\_statistical\_analysis\_outputs.py}. Two analysis scopes are kept separate:
\begin{enumerate}
    \item the full wav2vec2 synthetic sweep, containing 4,800 successful clip-level evaluations;
    \item the matched cross-model subset, containing 700 successful clips for each of the four
    ASR systems.
\end{enumerate}

This separation is necessary because wav2vec2 was evaluated on the full synthetic corpus, while
the other three systems have hypotheses only for the planned model-comparison subset. The
cross-model subset is the union of two controlled slices:
\begin{equation}
\begin{split}
\mathcal{S}_{\mathrm{model}} ={}&
\{\mathrm{SNR}=7.4,\ \mathrm{noise}=T,\ \mathrm{OVR}\in\{0.00,0.14,0.20,0.40\}\}\\
&\cup
\{\mathrm{OVR}=0.14,\ \mathrm{noise}=T,\ \mathrm{SNR}\in\{\mathrm{clean},7.4,0,-5\}\}.
\end{split}
\end{equation}
The subset is therefore a matched ``plus'' design rather than a complete factorial grid. This
is suitable for fair model ranking and for examining two one-factor profiles, but it should not
be interpreted as evidence about every SNR--overlap--noise combination for every model.

\section{Variable Selection and Justification}
\subsection{Dependent Variable}
The principal dependent variable is clip-level word error rate (WER). WER is selected as the
primary response variable for three reasons. First, it is the standard reporting metric in ASR
evaluation and is therefore interpretable against existing work. Second, it is available for all
four systems in the matched subset, including both segmented and non-segmented hypotheses. Third,
it has a direct operational interpretation: higher WER corresponds to more manual correction
effort for a fixed transcript length.

DSS-WER, ORC-WER and cpWER are retained as sensitivity checks rather than used as the primary
cross-model score. This is a conservative choice. The alternative metrics are valuable for
multi-speaker evaluation, but they also encode stronger assumptions about segmentation,
speaker-attribution, or reference serialisation. Ordinary WER is therefore the most stable
common denominator for the model-ranking analysis, while DSS-WER is used to test whether the
standard serialisation assumption materially changes the conclusions.

\subsection{Explanatory Variables}
The explanatory variables are signal-to-noise ratio (SNR), overlap ratio (OVR), noise type, and
model identity. SNR is included because additive acoustic interference is a direct stressor for
ASR. OVR is included because overlapping speech is the central phenomenon studied in this
dissertation. Noise type is included because two conditions with the same nominal SNR may differ
in spectral and temporal structure. Model identity is included because the evaluated systems
represent different architectural and decoding choices.

Speaker count is held at two in the synthetic experiment. This makes OVR interpretable as a
controlled measure of two-speaker concurrency rather than a mixture of overlap and speaker-count
effects. Clip duration is approximately \SI{60}{\second}, which supports comparable timing and
error estimates across conditions.

\section{Statistical Procedure}
All reported results use the clip as the unit of analysis. The design contains repeated
structure: the same base material is reused across acoustic variants, and the cross-model subset
uses the same clip IDs across systems. For this reason, the analysis emphasises matched
condition summaries, distributional plots and practical effect sizes rather than treating every
row as an independent sample from an unrelated population.

WER is non-negative and right-skewed, so means are reported together with medians and standard
deviations. Means are useful for estimating expected correction burden, while medians describe
the typical clip and are less sensitive to difficult outliers. The full-factor wav2vec2 analysis
is used for attribution of SNR, OVR and noise-type effects; cross-model claims are restricted to
the matched 700-clip subset.

\section{Full-Factor wav2vec2 Analysis}
\subsection{Effect of SNR}
Across the 4,800 successful wav2vec2 evaluations, the mean WER is 0.386 and the median WER is
0.320. The SNR effect is monotonic: mean WER rises from 0.224 in the clean condition to 0.588
at \(-5\) dB. This is an absolute increase of 0.364 WER, or approximately 2.6 times the clean
condition mean.

\begin{table}[htbp]
\centering
\caption{wav2vec2 WER by SNR on the full synthetic sweep.}
\label{tab:wav2vec2_snr_summary}
\begin{tabular}{lrrrr}
\toprule
\textbf{SNR condition} & \textbf{Clips} & \textbf{Mean WER} & \textbf{Median WER} & \textbf{SD} \\
\midrule
clean & 1,200 & 0.224 & 0.216 & 0.147 \\
7.4 dB & 1,200 & 0.279 & 0.266 & 0.167 \\
0 dB & 1,200 & 0.452 & 0.426 & 0.266 \\
\(-5\) dB & 1,200 & 0.588 & 0.576 & 0.313 \\
\bottomrule
\end{tabular}
\end{table}

\begin{figure}[htbp]
    \centering
    \includefigure[width=0.82\linewidth]{figures/stat_wav2vec2_snr_box.png}
    \caption{Distribution of clip-level WER for wav2vec2 by SNR condition.}
    \label{fig:stat_wav2vec2_snr_box}
\end{figure}

\Cref{fig:stat_wav2vec2_snr_box} shows that the degradation is not merely a shift in a small
number of extreme clips. Both the median and the upper tail move upward as the acoustic
condition worsens. The standard deviation also increases from 0.147 in clean audio to 0.313 at
\(-5\) dB, indicating that low-SNR transcription is less predictable as well as less accurate.

\subsection{Effect of Overlap and Noise Type}
The overlap effect is also monotonic in the wav2vec2 sweep. Mean WER rises from 0.246 at
OVR \(=0.00\) to 0.540 at OVR \(=0.40\). The corresponding medians are 0.071 and 0.473. The large
difference between mean and median at OVR \(=0.00\) reflects the fact that the zero-overlap
condition still includes noisy variants; non-overlap is not the same as clean audio.

\begin{table}[htbp]
\centering
\caption{wav2vec2 WER by overlap ratio on the full synthetic sweep.}
\label{tab:wav2vec2_ovr_summary}
\begin{tabular}{lrrrr}
\toprule
\textbf{OVR} & \textbf{Clips} & \textbf{Mean WER} & \textbf{Median WER} & \textbf{SD} \\
\midrule
0.00 & 1,200 & 0.246 & 0.071 & 0.317 \\
0.14 & 1,200 & 0.350 & 0.222 & 0.261 \\
0.20 & 1,200 & 0.407 & 0.296 & 0.234 \\
0.40 & 1,200 & 0.540 & 0.473 & 0.179 \\
\bottomrule
\end{tabular}
\end{table}

\begin{figure}[htbp]
    \centering
    \includefigure[width=0.95\linewidth]{figures/stat_wav2vec2_snr_overlap_heatmap.png}
    \caption{wav2vec2 mean WER and cell count across the SNR--OVR grid.}
    \label{fig:stat_wav2vec2_heatmap}
\end{figure}

\Cref{fig:stat_wav2vec2_heatmap} gives the clearest view of the joint response surface. The
lowest mean WER occurs for clean, non-overlapped speech (0.036), while the highest occurs at
OVR \(=0.40\) and \(-5\) dB (0.680). The sample-count panel confirms balanced support: each
SNR--OVR cell contains 300 clips. The interaction is not a simple linear amplification. At
severe noise levels, WER is already high even without overlap, so the additional overlap penalty
partly saturates. The practical conclusion is nevertheless clear: both lower SNR and higher OVR
move the model towards a higher-error operating region.

Noise type is a further source of variation. Public noise has the highest mean WER (0.569),
followed by domestic noise (0.328) and transportation noise (0.260). Since the cross-model
subset uses transportation noise only, the cross-model results should be read as a controlled
comparison slice rather than as an average over all noise environments.

\begin{figure}[htbp]
    \centering
    \includefigure[width=0.68\linewidth]{figures/stat_wav2vec2_noise_type.png}
    \caption{wav2vec2 mean WER by DEMAND noise category. Error bars show approximate 95\% confidence intervals.}
    \label{fig:stat_wav2vec2_noise}
\end{figure}

\section{Matched Cross-Model Comparison}
\subsection{Pooled WER Ranking}
The matched model-comparison subset contains 700 clips per model. The same clip IDs are used for
all four systems, so differences in the pooled table reflect model behaviour on a common
evaluation set rather than different data composition.

% \begin{table}[htbp]
% \centering
% \caption{WER-first model comparison on the matched 700-clip subset.}
% \label{tab:cross_model_wer_first}
% \begin{tabular}{lrrrrr}
% \toprule
% \textbf{Model} & \textbf{Clips} & \textbf{Mean WER} & \textbf{Median WER} & \textbf{SD} & \textbf{\% WER \(>1\)} \\
% \midrule
% faster-whisper & 700 & 0.159 & 0.123 & 0.128 & 0.14 \\
% WhisperX & 700 & 0.181 & 0.139 & 0.150 & 0.00 \\
% wav2vec2 & 700 & 0.228 & 0.205 & 0.133 & 0.00 \\
% Parakeet & 700 & 0.503 & 0.525 & 0.208 & 0.00 \\
% \bottomrule
% \end{tabular}
% \end{table}

\begin{figure}[htbp]
    \centering
    \includefigure[width=0.82\linewidth]{figures/stat_cross_model_wer_box.png}
    \caption{Distribution of WER by model on the matched cross-model subset.}
    \label{fig:stat_cross_model_box}
\end{figure}

\Cref{tab:cross_model_wer_first} and \Cref{fig:stat_cross_model_box} show that faster-whisper
has the lowest mean and median WER on this subset, followed by WhisperX, wav2vec2 and Parakeet.
The gap between faster-whisper and WhisperX is modest (0.023 mean WER), while the gap between
the first three systems and Parakeet is much larger. The model ranking is therefore driven by
broad distributional shifts rather than by a small number of catastrophic cases.

\subsection{Condition Profiles Within the Matched Subset}
The pooled ranking hides two different controlled profiles. The first profile varies OVR while
holding SNR at 7.4 dB and noise type at transportation. The second varies SNR while holding OVR
at 0.14 and noise type at transportation.

\begin{figure}[htbp]
    \centering
    \includefigure[width=0.96\linewidth]{figures/stat_cross_model_condition_profiles.png}
    \caption{Cross-model WER profiles for the two matched comparison slices.}
    \label{fig:stat_cross_model_profiles}
\end{figure}

In the overlap slice, the three lower-WER systems show the expected upward trend as OVR
increases. At OVR \(=0.40\), faster-whisper, WhisperX and wav2vec2 have mean WERs of 0.339,
0.383 and 0.431 respectively. Parakeet remains substantially higher across the slice, with mean
WER 0.598 at OVR \(=0.40\). In the SNR slice, faster-whisper and WhisperX are comparatively
stable across the four SNR levels, wav2vec2 degrades more clearly as SNR falls, and Parakeet
shows a non-monotonic pattern. This exception is a useful warning against over-generalising from
the pooled cross-model subset: the matched subset supports fair ranking under the selected
conditions, but wav2vec2 remains the stronger basis for full-factor acoustic attribution.

\section{Metric Sensitivity}
The primary analysis uses WER, but the metric-sensitivity check compares WER against DSS-WER on
the same 700 matched clips. In the JSON outputs, DSS-WER is stored as \texttt{mrs\_wer}. Define
\begin{equation}
    d_i = \mathrm{DSS\mbox{-}WER}_i - \mathrm{WER}_i.
\end{equation}
A negative value of \(d_i\) means that DSS-WER is more forgiving than ordinary WER for that clip.

\begin{table}[htbp]
\centering
\caption{Metric sensitivity on the matched cross-model subset.}
\label{tab:metric_sensitivity}
\begin{tabular}{lrrr}
\toprule
\textbf{Model} & \textbf{Mean WER} & \textbf{Mean DSS-WER} & \textbf{Mean difference \(d\)} \\
\midrule
faster-whisper & 0.159 & 0.156 & -0.003 \\
WhisperX & 0.181 & 0.179 & -0.002 \\
wav2vec2 & 0.228 & 0.211 & -0.017 \\
Parakeet & 0.503 & 0.515 & 0.011 \\
\bottomrule
\end{tabular}
\end{table}

\Cref{tab:metric_sensitivity} shows that DSS-WER does not overturn the main WER ranking. For
faster-whisper and WhisperX, the mean difference is very small. For wav2vec2, DSS-WER is lower
by 0.017 on average, suggesting that some of its ordinary WER penalty is attributable to
serialisation effects. Parakeet is the exception: its DSS-WER mean is slightly higher than its
ordinary WER mean on this subset. This supports a cautious interpretation of the proposed
metric. DSS-WER is useful as a diagnostic for overlapped references, but it should be reported
alongside ordinary WER rather than used as a replacement ranking metric.

\section{Efficiency Analysis}
Accuracy is the primary criterion, but inference cost matters for deployment. The timing
benchmark uses 100 randomly sampled clips and reports both time per clip and a real-time
multiple:
\begin{equation}
    s_m = \frac{T_m}{N_m},
    \qquad
    \rho_m = \frac{D_m}{T_m},
\end{equation}
where \(T_m\) is total processing time, \(N_m\) is the number of clips and \(D_m\) is aggregate
audio duration. Larger \(\rho_m\) indicates higher throughput.

% \begin{table}[htbp]
% \centering
% \caption{Inference efficiency on the 100-clip timing benchmark.}
% \label{tab:timing_efficiency}
% \begin{tabular}{lrr}
% \toprule
% \textbf{Model} & \textbf{Seconds per clip} & \textbf{Real-time multiple \(\rho\)} \\
% \midrule
% wav2vec2 & 0.754 & 88.951 \\
% Parakeet & 1.124 & 59.663 \\
% WhisperX & 7.276 & 9.220 \\
% faster-whisper & 27.736 & 2.419 \\
% \bottomrule
% \end{tabular}
% \end{table}

\begin{figure}[htbp]
    \centering
    \includefigure[width=0.72\linewidth]{figures/stat_accuracy_latency_tradeoff.png}
    \caption{Accuracy-latency trade-off using matched-subset WER and the 100-clip timing benchmark.}
    \label{fig:stat_accuracy_latency}
\end{figure}

\Cref{tab:timing_efficiency} and \Cref{fig:stat_accuracy_latency} show a clear trade-off.
wav2vec2 is the fastest system by a large margin, but it is not the most accurate on the matched
subset. faster-whisper has the lowest WER but the highest latency. WhisperX lies between these
two systems. Parakeet is fast, but its WER on the matched subset is high enough that speed alone
does not make it the preferred system for this evaluation setting.

\section{Real-Audio External Validity Check}
The CHiME-6 real-audio evaluation contains 100 clips per model in
\texttt{Output/real\_eval\_results.json}. These results are included only as an external
validity check, because the scoring metric differs by output structure: segmented outputs are
reported with cpWER, while single-stream outputs are reported with ordinary WER.

\begin{table}[htbp]
\centering
\caption{Real-audio scores on the CHiME-6 segment set. Metrics differ by output structure and are not a strict WER-only ranking.}
\label{tab:real_audio_summary}
\begin{tabular}{llrr}
\toprule
\textbf{Model} & \textbf{Metric} & \textbf{Clips} & \textbf{Mean score} \\
\midrule
Parakeet & WER & 100 & 0.389 \\
wav2vec2 & WER & 100 & 0.723 \\
WhisperX & cpWER & 100 & 0.836 \\
faster-whisper & cpWER & 100 & 0.895 \\
\bottomrule
\end{tabular}
\end{table}

The real-audio results do not reproduce the synthetic matched-subset ranking. This is not a
contradiction of the synthetic analysis; rather, it indicates limited transfer from controlled
mixtures to far-field conversational audio. Real meetings include channel effects, natural turn
taking, diarisation ambiguity and reference-alignment challenges that are not fully represented
by the synthetic factorial design. The defensible conclusion is therefore two-stage: synthetic
sweeps are appropriate for controlled factor attribution, while final deployment decisions
require validation on real meeting audio.

\section{Limitations}
Three limitations should be made explicit. First, the cross-model subset is matched but not a
complete factorial grid, so full SNR--OVR--noise interactions are estimated only for wav2vec2.
Second, the model-comparison subset uses transportation noise, which is the easiest noise type
for wav2vec2 in the full sweep; cross-model WERs should therefore not be treated as averages over
all acoustic environments. Third, ordinary WER is necessary for comparability but incomplete for
overlapped speech, which is why DSS-WER is reported as a sensitivity analysis rather than
ignored.

\section{Conclusion}
The statistical evidence supports the central robustness claim of the dissertation. In the full
wav2vec2 sweep, lower SNR, higher overlap ratio and noise type all materially affect WER. In the
matched cross-model subset, faster-whisper gives the lowest WER, WhisperX is close behind,
wav2vec2 trades accuracy for speed, and Parakeet is fast but substantially less accurate under
the selected synthetic conditions. The combined analysis justifies evaluating ASR robustness as
a joint function of acoustic condition, model identity, metric policy and computational cost,
rather than relying on a single pooled WER value.
